\chapter*{Preface}
\addcontentsline{toc}{chapter}{Preface}

Most books about programming languages fall into one of two camps, and neither one quite helps a working engineer make the jump from an older object-oriented language to modern systems and type-driven software.

The first camp is the syntax tour: ``here is how you write a loop in Rust, here is how you declare a class in Kotlin, here is how you spawn a goroutine in Go.'' After fifty pages you can transcribe small programs, yet the first time the Rust borrow checker rejects a perfectly ordinary listener callback, or TypeScript narrows a union type in a way Java generics cannot express, or a Go interface holding a \texttt{nil} pointer evaluates \texttt{err != nil} to \texttt{true}, you hit a wall. Syntax was never the hard part. The mental model underneath the syntax was.

The second camp is pure programming-language theory: operational semantics, lambda calculus reductions, and inference rules stacked like fractions before you ever inspect a single byte in memory or send a packet over a socket. Theory matters --- algebraic data types, parametric polymorphism, variance, and linear ownership are clean mathematical ideas --- yet studying those ideas detached from stack frames, CPU cache lines, vtable pointers, and real production failure modes strips away the engineering pressure that forced language designers to invent them.

This book works in both layers at once. It keeps one eye on the hardware --- exact stack frames, 64-byte cache lines, object headers, fat pointers, state-machine layouts --- and the other on the type system and mathematical structure that let a compiler rule out whole categories of bugs before an executable exists. The running example throughout is a concrete production system: a high-assurance financial ledger backend in Rust with Axum, SQLx, and PostgreSQL.

\section*{Who this book is for, and where we start}

We assume a very specific, realistic starting point:
\begin{itemize}
  \item \textbf{You know Java at roughly the Java~6 level.} You are comfortable with classes, interfaces, inheritance (\texttt{extends}, \texttt{implements}), \texttt{ArrayList} and \texttt{HashMap} generics, \texttt{null}, checked and unchecked \texttt{try}/\texttt{catch}/\texttt{finally} exceptions, \texttt{synchronized} blocks, OS threads, and garbage collection. You do \emph{not} need Java~8 streams, records, or sealed classes --- when modern JVM ideas appear, we build them cleanly from Java~6 upward.
  \item \textbf{You know basic C and basic C++.} You have written \texttt{for} loops, \texttt{if}/\texttt{else} branches, functions, \texttt{struct} declarations, arrays, \texttt{malloc}/\texttt{free}, and raw pointers (\texttt{T*} and \texttt{\&x}). You do \emph{not} need modern C++11--C++20 move semantics, smart pointers, templates, or RAII --- Chapter~2 builds all of those from scratch.
\end{itemize}

Because you already know how loops, conditionals, and functions work, we spend zero pages re-explaining beginner control flow. Every page goes toward the questions that actually rewire how you read and write software:

\begin{itemize}
  \item When you write a variable declaration in Java~6, C++, Rust, Go, Kotlin, or TypeScript, \emph{what exact bytes sit on the stack, what sits on the heap, how many pointer dereferences stand between the CPU and your integer, and who frees those bytes when?}
  \item Why did Rust reject class inheritance, \texttt{null}, exceptions, and garbage collection --- and how do \emph{ownership}, \emph{borrowing} (\texttt{\&T} versus \texttt{\&mut T}), \emph{lifetimes} (\texttt{'a}), \emph{algebraic data types} (\texttt{struct} and \texttt{enum}), and \emph{traits} fit together into a single coherent lens?
  \item How do functional programmers in \textbf{OCaml} and \textbf{Haskell} actually think? Once you strip away academic jargon, how do immutability, algebraic variants, pattern matching, higher-order pipelines, typeclasses, functors, and monads show up directly inside everyday idiomatic Rust?
  \item How do \textbf{TypeScript} (structural set-theoretic types erased at runtime), \textbf{Kotlin} (null-safe algebraic hierarchies and compiler-rewritten CPS coroutines on the JVM), and \textbf{Go} (flat value structs, escape analysis, implicit fat-pointer interfaces, and an M:N goroutine scheduler) each solve the exact pain points of Java~6 with radically different engineering tradeoffs?
  \item How do all of these low-level and type-level tools pay off when real money is on the line --- building an asynchronous \textbf{Rust + Axum + SQLx + PostgreSQL} double-entry financial ledger where negative transfers, currency mismatches, lost concurrent updates, unbalanced journal entries, and unrolled-back partial failures are prevented by construction?
\end{itemize}

\section*{Why a financial ledger is our north star}

Toy examples flatter weak abstractions. Almost any language looks pleasant when you are printing \texttt{Animal.speak()} for a \texttt{Dog} and a \texttt{Cat}. Financial systems are unforgiving in a much more instructive way:

\begin{enumerate}
  \item \textbf{Representation precision matters.} Using IEEE-754 binary floating-point (\texttt{double} / \texttt{f64}) for money produces $0.10 + 0.20 = 0.30000000000000004$. Silent integer overflow turns a large credit into a large negative balance. Mixing USD cents with EUR cents or swapping an \texttt{AccountId} with a \texttt{JournalId} corrupts books silently unless types keep them distinct at zero runtime cost.
  \item \textbf{Illegal states must be unrepresentable.} In a Java~6 \texttt{PaymentOrder} class with eight nullable fields and three boolean flags (\texttt{isVerified}, \texttt{isSettled}, \texttt{isFailed}), $2^3 = 8$ flag combinations exist even when only three make business sense, and every method starts with defensive \texttt{if (order.getSettledAt() != null)} checks. Sum types (\texttt{enum}s) and typestate collapse invalid combinations at compile time.
  \item \textbf{Concurrency and partial failure are adversarial.} If two HTTP requests simultaneously withdraw $\$80$ from an account holding $\$100$, or if a client drops its TCP connection halfway through an \texttt{async} transfer handler, what happens to the lock, the memory, and the uncommitted PostgreSQL transaction? Answering that requires understanding stackless \texttt{Future} state machines, \texttt{Drop}/RAII guards, and database isolation together.
\end{enumerate}

Even in Parts~I through~IV, before we reach the full service in Part~V, you will meet orders, accounts, minor-unit money amounts, and ledger postings alongside systems classics like buffers, parsers, and schedulers, so that by Chapter~12 the capstone uses only ideas you have already met.

\section*{How to read this book, and what the boxes mean}

Every concept in this book follows a steady four-step rhythm: \emph{the concrete bug or hardware cost first}, \emph{the physical byte-level picture second}, \emph{the language or type-system rule third}, and \emph{idiomatic code side by side across languages fourth}.

Throughout the chapters, recurring colored boxes tell you what kind of thinking a passage asks for:

\begin{intuitionbox}
The core mental picture --- the physical or geometric way of seeing a rule so you never have to memorize it by rote. If you revisit a chapter months later, reading the intuition boxes reconstructs the whole arc.
\end{intuitionbox}

\begin{focusbox}
A central invariant or guiding principle that organizes an entire section or language design choice --- for instance, Rust's \emph{Aliasing XOR Mutability} rule or TypeScript's \emph{types as sets of values}.
\end{focusbox}

\begin{formalbox}
The exact rule, memory layout, or type relation stated without hand-waving --- byte offsets, alignment rules, subtyping judgments, or compiler desugarings. When intuition and edge cases collide, this box gives the ground truth.
\end{formalbox}

\begin{whybox}
The engineering and mathematical reason \emph{why} a rule is designed the way it is --- why C++ move constructors cannot simply skip the moved-from destructor, or why mutable references in Rust and mutable generic containers in Java/TypeScript must be invariant.
\end{whybox}

\begin{trapbox}
A classic footgun that catches smart engineers moving from Java~6 or C/C++ into modern languages --- shown with the exact code that triggers it and the mechanism behind the failure.
\end{trapbox}

\begin{labbox}
A concrete trace or micro-experiment to work through on paper or at a terminal: counting cache misses, writing down a compiler-generated closure struct, or hand-stepping an \texttt{async} state machine across \texttt{.await} points.
\end{labbox}

\begin{applicationbox}
How the mechanism behaves in real production systems: tail-latency spikes in trading and payment engines, database connection pools, serialization boundaries, and network servers.
\end{applicationbox}

You will also occasionally see \emph{History and connection} boxes (tracing how ideas traveled from ML, Haskell, and C++ into Rust, Kotlin, TypeScript, and Go), \emph{Under the hood} boxes (peeking at assembly, JVM bytecode, or desugared compiler output), \emph{Walking question} boxes (a puzzle to turn over in your head), and \emph{Pause} boxes (checkpoints summarizing what you can now do).

\section*{A ten-minute preview: one transfer, three mental models}

To taste the difference between reading syntax and reading through the machine-and-types lens, look at the same miniature operation --- debiting an account balance --- written in Java~6, basic C, and Rust.

First, a familiar Java~6 method:

\begin{lstlisting}[language=Java]
// Java 6: heap reference, implicit null, exception control flow
public class Account {
    private Long balanceCents; // boxed heap object!

    public void debit(Long amountCents) throws InsufficientFundsException {
        if (amountCents <= 0) {
            throw new IllegalArgumentException("amount must be positive");
        }
        if (this.balanceCents < amountCents) {
            throw new InsufficientFundsException(this.balanceCents, amountCents);
        }
        this.balanceCents = this.balanceCents - amountCents;
    }
}
\end{lstlisting}

Look past the clean surface and ask what the JVM actually executes, and what the type signature hides:
\begin{enumerate}
  \item \textbf{Hidden \texttt{null}s.} Both \texttt{balanceCents} and \texttt{amountCents} are boxed \texttt{java.lang.Long} references. If a caller passes \texttt{null}, the check \texttt{amountCents <= 0} calls \texttt{amountCents.\allowbreak longValue()} under the hood and throws \texttt{NullPointerException} --- an exit path mentioned nowhere in the signature.
  \item \textbf{Pointer chasing and allocation.} Reading \texttt{this.balanceCents} follows a pointer to the \texttt{Account} object (with its 16-byte JVM object header), then follows \emph{another} pointer to a separate heap-allocated \texttt{Long} object (another 16-byte header plus 8 payload bytes = 24 bytes just to hold an 8-byte integer). Subtracting \texttt{amountCents} allocates a brand-new \texttt{Long} object on the heap unless the result happens to fall in \texttt{Long.valueOf}'s tiny $[-128, 127]$ cache.
  \item \textbf{Unrestricted aliasing.} Any number of threads, lists, and caches may hold references to this same \texttt{Account} simultaneously while \texttt{debit} mutates it in place. Nothing in the type tells you whether you hold the only handle or one of fifty.
\end{enumerate}

Now swing to the opposite extreme in C:

\begin{lstlisting}[language=C]
/* C: flat bytes in memory, zero GC, zero compile-time guardrails */
#include <stdint.h>

typedef struct {
    uint64_t id;
    int64_t  balance_cents;
} Account;

int account_debit(Account *acct, int64_t amount_cents) {
    if (acct == NULL || amount_cents <= 0) return -1;
    if (acct->balance_cents < amount_cents) return -2;
    acct->balance_cents -= amount_cents;
    return 0;
}
\end{lstlisting}

The hardware picture is now plain: \texttt{Account} is a flat 16-byte value (\texttt{id} at offset 0, \texttt{balance\_cents} at offset 8), with zero object headers and zero heap allocations during subtraction. Yet notice what we gave up to get that flatness:
\begin{itemize}
  \item \texttt{acct} might still be \texttt{NULL}, or worse, a \emph{dangling pointer} to a stack frame that already returned or a heap block already passed to \texttt{free()}.
  \item The error codes \texttt{-1} and \texttt{-2} sit in an ordinary \texttt{int}. A caller who writes \texttt{account\_debit(\&a, 5000);} and forgets to check the return value continues executing as if the debit succeeded.
  \item Nothing stops two pointers from aliasing the same \texttt{Account} across threads or across function arguments while mutating it.
\end{itemize}

Now look at the idiomatic Rust version we will be able to read down to the register level by the end of Part~II:

\begin{lstlisting}[language=Rust]
// Rust: flat C-like memory layout + compile-time ownership & algebraic types
#[derive(Debug, Clone, Copy, PartialEq, Eq)]
pub struct PositiveCents(i64); // invariant: always > 0

#[derive(Debug, PartialEq, Eq)]
pub enum DebitError {
    InsufficientFunds { available: i64, requested: i64 },
    ArithmeticOverflow,
}

pub struct Account {
    pub id: u64,
    balance_cents: i64,
}

impl Account {
    pub fn debit(&mut self, amount: PositiveCents) -> Result<(), DebitError> {
        if self.balance_cents < amount.0 {
            return Err(DebitError::InsufficientFunds {
                available: self.balance_cents,
                requested: amount.0,
            });
        }
        self.balance_cents = self.balance_cents
            .checked_sub(amount.0)
            .ok_or(DebitError::ArithmeticOverflow)?;
        Ok(())
    }
}
\end{lstlisting}

Pause and count how many runtime checks and heap allocations disappeared \emph{at the same time}:
\begin{itemize}
  \item \textbf{Exact C memory layout.} \texttt{Account} occupies the exact same 16 flat bytes as the C \texttt{struct}, and \texttt{PositiveCents(i64)} compiles down to a bare 64-bit register value with zero wrapper header.
  \item \textbf{No \texttt{null}, no dangling pointer, no aliasing race.} The parameter \texttt{\&mut self} is a compile-time proof that \texttt{self} points to a live, properly aligned \texttt{Account} \emph{and} that no other active reference anywhere in the program can read or write this \texttt{Account} for the duration of \texttt{debit}.
  \item \textbf{Invariants baked into types.} Because \texttt{PositiveCents}'s inner \texttt{i64} is private to its module and can only be constructed through a smart constructor that checks \texttt{val > 0} once at the boundary, \texttt{debit} never even needs an \texttt{amount <= 0} check.
  \item \textbf{Errors as flat values.} \texttt{Result<(), DebitError>} is a stack-allocated tagged union, not a heap-allocated exception unwinding stack frames. And if a caller ignores a \texttt{Result}, the Rust compiler immediately warns that a \texttt{\#[must\_use]} failure value was dropped.
\end{itemize}

\begin{center}
\begin{tikzpicture}[
    font=\small,
    >=Stealth,
    colbox/.style={draw=#1!80!black, fill=#1!8, rounded corners=3pt, thick, align=left, inner sep=6pt, text width=3.95cm, minimum height=3.55cm}
  ]
  \node[colbox=quotientorange] (java) at (0, 0) {
    \textbf{\footnotesize\color{quotientorange!85!black}Java 6 (\texttt{Account} + \texttt{Long})}\\[3pt]
    \scriptsize
    $\bullet$ Heap object + 16B header\\
    $\bullet$ Boxed \texttt{Long}: 2nd heap hop\\
    $\bullet$ Hidden \texttt{null} + exceptions\\
    $\bullet$ Shared mutable aliases\\
    $\bullet$ Safe memory via GC pauses
  };

  \node[colbox=trapred] (clang) at (4.65, 0) {
    \textbf{\footnotesize\color{trapred!85!black}C (\texttt{Account*})}\\[3pt]
    \scriptsize
    $\bullet$ Flat 16B \texttt{struct}, 0 headers\\
    $\bullet$ Direct register arithmetic\\
    $\bullet$ Raw \texttt{NULL} \& dangling ptrs\\
    $\bullet$ Ignorable \texttt{int} error codes\\
    $\bullet$ Zero memory safety
  };

  \node[colbox=targetgreen] (rust) at (9.3, 0) {
    \textbf{\footnotesize\color{targetgreen!85!black}Rust (\texttt{\&mut Account})}\\[3pt]
    \scriptsize
    $\bullet$ Flat 16B \texttt{struct}, 0 headers\\
    $\bullet$ \texttt{PositiveCents} = 8B register\\
    $\bullet$ \texttt{\&mut}: live + unique access\\
    $\bullet$ \texttt{Result} stack sum type\\
    $\bullet$ Compile-time safety, 0 GC
  };

  \draw[->, thick, sourceblue!80!black] (java.north) to[bend left=22]
    node[above, font=\scriptsize\bfseries, text=sourceblue!85!black] {drop headers \& GC (Part I)} (clang.north);
  \draw[->, thick, targetgreen!80!black] (clang.south) to[bend right=22]
    node[below, font=\scriptsize\bfseries, text=targetgreen!85!black] {add ownership \& sum types (Part II)} (rust.south);
\end{tikzpicture}
\end{center}

Seeing \emph{both} layers at once --- the flat 16 bytes on the machine and the mathematical guarantees in the type signature --- is what we mean by ``from the ground up.''

\section*{Roadmap of the five parts}

The fourteen chapters of this book are grouped into five progressive parts:

\textbf{Part~I --- Leaving Java~6: Memory, Values, and the Machine (Chapters~1--2).}
Chapter~1 strips away the JVM illusion and inspects what Java~6 actually puts in RAM: stack frames versus heap objects, 16-byte object headers, pointer chasing across 64-byte CPU cache lines, stop-the-world GC pauses, aliasing bugs, \texttt{null}, and exception control flow, contrasted directly against flat value arrays in C, C++, and Rust. Chapter~2 builds \textbf{Modern C++ (C++11--C++20)} from basic C/C++: deterministic RAII destructors, lvalue/rvalue references, copy versus move semantics (\texttt{std::move} and why C++ moves are \emph{non-destructive}), \texttt{unique\_ptr}/\texttt{shared\_ptr}/\texttt{string\_view}/\texttt{span}, monomorphized templates with C++20 Concepts, lambdas, and the exact four C++ footguns (dangling views, iterator invalidation on \texttt{vector} reallocation, use-after-move, and data races) that Rust was invented to eliminate at compile time.

\textbf{Part~II --- Seeing Through the Rust Lens (Chapters~3--6).}
Chapter~3 derives Rust's single foundational invariant --- \emph{Aliasing XOR Mutability} --- and maps out the exact stack and heap byte layouts of \texttt{String}, \texttt{\&str}, \texttt{Vec<T>}, \texttt{\&[T]}, \texttt{Box<T>}, \texttt{Rc<T>}, \texttt{Arc<T>}, \texttt{Cell<T>}, \texttt{RefCell<T>}, and \texttt{Mutex<T>}. Chapter~4 removes all mystery from lifetimes (\texttt{'a}), Non-Lexical Lifetimes, elision, \texttt{'static}, variance, and the \texttt{unsafe} proof boundary. Chapter~5 develops algebraic data types --- product types (\texttt{struct}) and sum types (\texttt{enum}), niche layout optimizations where \texttt{Option<\&T>} occupies just 8 bytes, \texttt{Option}/\texttt{Result}/\texttt{?}, and exhaustive pattern matching. Chapter~6 opens up traits, generics, static dispatch via monomorphization versus dynamic dispatch via 16-byte \texttt{(data\_ptr, vtable\_ptr)} fat pointers, coherence rules, and standard library traits.

\textbf{Part~III --- Thinking Functionally and Writing Idiomatic Rust (Chapters~7--8).}
Chapter~7 rewires imperative loop-and-flag habits by visiting Rust's direct functional ancestors: \textbf{OCaml} (immutability by default, expression-oriented control flow, Hindley--Milner type inference, algebraic variants, persistent structural sharing, higher-order pipelines) and \textbf{Haskell} (purity, laziness without thunks in Rust iterators, Typeclasses as the direct parent of traits, Functors as \texttt{.map()}, and Monads as \texttt{.and\_then()} and \texttt{?}), ending with how Rust compiles closures (\texttt{Fn}, \texttt{FnMut}, \texttt{FnOnce}) and iterator chains into single zero-allocation machine loops. Chapter~8 shows how experienced engineers actually architect Rust systems: curing ``Java-in-Rust'' cyclic OOP graphs with flat \texttt{Vec} arenas and typed index handles, \emph{Parse, Don't Validate}, newtypes, builders, zero-byte compile-time typestate (\texttt{PhantomData}), borrowing views at API boundaries, and \texttt{thiserror}/\texttt{anyhow} error design.

\textbf{Part~IV --- Modern Type Systems and Concurrency: TypeScript, Kotlin, and Go (Chapters~9--11).}
Once you hold both the hardware lens and the Rust/FP lens, other modern languages snap into focus rapidly. Chapter~9 builds \textbf{TypeScript} from set theory and structural subtyping: \texttt{never} ($\emptyset$), \texttt{unknown} ($U$), unions ($A \cup B$), intersections ($A \cap B$), discriminated unions, control-flow narrowing, mapped and conditional types (\texttt{keyof}, \texttt{infer}), branded types, and runtime validation (\texttt{Zod}) across the erase-to-JS boundary. Chapter~10 bridges Java~6 directly to \textbf{Kotlin} down to JVM bytecode: nullability in the type system, \texttt{@JvmInline value class}, \texttt{sealed interface}, \texttt{inline}/\texttt{reified} generics defeating JVM type erasure, and how \texttt{suspend} functions compile via Continuation-Passing Style (CPS) into state machines under structured concurrency. Chapter~11 dissects \textbf{Go} (value structs, escape analysis, fat-pointer interfaces and the typed-\texttt{nil} trap, and the M:N \texttt{G-M-P} goroutine scheduler with channels) and closes with a seven-axis architectural synthesis across Java~6, Modern C++, Rust, OCaml, Haskell, TypeScript, Kotlin, and Go.

\textbf{Part~V --- Capstone: A Financial Backend in Rust with Axum, SQLx, and PostgreSQL (Chapters~12--14).}
Chapter~12 goes under the hood of \textbf{Async Rust and Tokio}: how \texttt{async fn} compiles into a stackless \texttt{enum} state machine, why self-referential borrows across \texttt{.await} require \texttt{Pin<\&mut Self>}, work-stealing scheduling, \texttt{Send}/\texttt{'static} bounds, backpressure, and cancellation safety when clients disconnect mid-request. Chapter~13 models \textbf{money and double-entry ledgers} in types: minor-unit integers and fixed-point decimals, checked arithmetic, banker's rounding with remainder allocation, currency safety, conservation of debits and credits ($\sum \text{debits} = \sum \text{credits}$), and typestate proof tokens (\texttt{BalancedEntry}). Chapter~14 assembles the complete \textbf{Axum + Tower + SQLx + PostgreSQL} financial service: production DDL, compile-time SQL type-checking (\texttt{sqlx::query!}), RAII transaction rollback on \texttt{Drop}, defeating concurrent overdraft races via canonically ordered \texttt{SELECT ... FOR UPDATE} and Stripe-style idempotency keys, transactional outbox delivery, graceful shutdown, and integration tests.

\begin{center}
\begin{tikzpicture}[
    font=\small,
    >=Stealth,
    pbox/.style={draw=#1!80!black, fill=#1!10, rounded corners=3pt, thick, align=center, inner sep=5pt, minimum height=1.2cm},
    arr/.style={->, thick, draw=gray!70!black}
  ]
  \node[pbox=sourceblue, minimum width=11.6cm] (p1) at (0, 0) {
    \textbf{\color{sourceblue!90!black}Part I: Leaving Java 6 --- Memory, Values \& Modern C++ (Ch.\ 1--2)}\\
    \scriptsize Stack vs.\ heap, 16B headers, cache lines, RAII, C++ moves vs.\ footguns
  };

  \node[pbox=targetgreen, minimum width=11.6cm] (p2) at (0, -1.75) {
    \textbf{\color{targetgreen!90!black}Part II: Seeing Through the Rust Lens (Ch.\ 3--6)}\\
    \scriptsize Ownership, \texttt{\&T}/\texttt{\&mut T}, lifetimes \texttt{'a}, \texttt{enum} sum types, traits \& fat pointers
  };

  \node[pbox=structurepurple, minimum width=5.55cm] (p3) at (-3.025, -3.65) {
    \textbf{\color{structurepurple!90!black}Part III: FP \& Idiomatic Rust}\\
    \textbf{\color{structurepurple!90!black}(Ch.\ 7--8)}\\
    \scriptsize OCaml, Haskell, closures, iterators,\\[-1pt]
    \scriptsize arenas, newtypes, typestate
  };

  \node[pbox=quotientorange, minimum width=5.55cm] (p4) at (3.025, -3.65) {
    \textbf{\color{quotientorange!90!black}Part IV: Modern Languages}\\
    \textbf{\color{quotientorange!90!black}(Ch.\ 9--11)}\\
    \scriptsize TypeScript set types, Kotlin CPS\\[-1pt]
    \scriptsize coroutines, Go M:N runtime \& synthesis
  };

  \node[pbox=trapred, minimum width=11.6cm] (p5) at (0, -5.65) {
    \textbf{\color{trapred!90!black}Part V: Capstone Financial Backend in Rust (Ch.\ 12--14)}\\
    \scriptsize Async \texttt{Future}s \& \texttt{Pin}, Tokio, double-entry ledger types, Axum + SQLx + PostgreSQL
  };

  \draw[arr] (p1.south) -- (p2.north);
  \draw[arr] (p2.south) -- ++(0,-0.22) -| (p3.north);
  \draw[arr] (p2.south) -- ++(0,-0.22) -| (p4.north);
  \draw[arr] (p3.south) |- ([yshift=0.22cm]p5.north) -- (p5.north);
  \draw[arr] (p4.south) |- ([yshift=0.22cm]p5.north) -- (p5.north);
\end{tikzpicture}
\end{center}

\begin{connectionbox}
When Chapter~14 locks rows in PostgreSQL, evaluates isolation levels, and stores idempotency keys and double-entry postings, it meets the companion volume \emph{Databases and PostgreSQL From the Ground Up} head-on. Here we examine database interactions from the application language and type-system side; the database reader opens up PostgreSQL's 8~KB heap pages, MVCC tuple headers, write-ahead log, and query planner underneath those exact SQL statements.
\end{connectionbox}

\section*{How to get the most out of every chapter}

Do not read passively. Whenever a diagram shows a 24-byte \texttt{Vec<T>} header on the stack pointing to a heap buffer, grab a scrap of paper and trace what happens to every word when the vector reallocates or moves into another function. Whenever a chapter contrasts Java~6, C++, and Rust, ask yourself \emph{which bugs failed at runtime in the first program, which bugs became undefined behavior in the second, and which exact rule rejected them at compile time in the third}.

Every chapter closes with exercises that ask you to calculate byte sizes and alignments by hand, spot subtle aliasing or concurrency bugs, predict compiler errors before running a compiler, and refactor brittle imperative designs into clean type-checked state machines. Working through them is where fluency sticks. To set that habit right away, try these four short warm-ups using only what we just previewed.

\section*{Exercises}

\begin{exercise}[Counting hidden exit paths in Java~6]
Go back to the ten-line Java~6 \texttt{Account.debit} method in this preface. List every distinct way a call \texttt{acct.debit(amt)} can terminate \emph{without} returning normally --- including what happens if \texttt{acct} itself is \texttt{null}, if \texttt{amt} is \texttt{null}, or if \texttt{balanceCents} was left \texttt{null} by an incomplete constructor. Which of those exit paths are visible in the method's type signature?
\end{exercise}

\begin{exercise}[Silent overflow with signed 64-bit cents]
Suppose the Java~6 \texttt{Account} class used primitive \texttt{long} instead of \texttt{Long}, and a separate \texttt{credit(long amount)} method updated \texttt{balanceCents += amount} without a range check. In Java (and two's-complement hardware), what happens if \texttt{balanceCents} equals \texttt{Long.MAX\_VALUE} ($9{,}223{,}372{,}036{,}854{,}775{,}807$) and \texttt{credit(100)} executes? Contrast this silent wraparound with \texttt{checked\_add} returning \texttt{Option<i64>} in Rust.
\end{exercise}

\begin{exercise}[The self-transfer aliasing bug]
Consider this short C transfer function that reads both balances before writing:
\begin{lstlisting}[language=C]
void transfer(Account *from, Account *to, int64_t cents) {
    int64_t f = from->balance_cents;
    int64_t t = to->balance_cents;
    if (f >= cents) {
        from->balance_cents = f - cents;
        to->balance_cents   = t + cents;
    }
}
\end{lstlisting}
Trace line by line what happens when a caller invokes \texttt{transfer(\&alice, \&alice, 5000)} with \texttt{alice}'s starting balance equal to \texttt{10000}. What is \texttt{alice}'s final balance, and where did the extra $\$50.00$ come from? Why does a Rust signature taking \texttt{from: \&mut Account} and \texttt{to: \&mut Account} reject \texttt{transfer(\&mut alice, \&mut alice, ...)} at compile time?
\end{exercise}

\begin{exercise}[Floating-point money on a napkin]
In IEEE-754 double-precision binary floating point (\texttt{double} in Java/C, \texttt{f64} in Rust), the decimal fraction $0.1 = \frac{1}{10}$ has the infinite repeating binary expansion $0.0001100110011\dots_2$, just as $\frac{1}{3} = 0.3333\dots_{10}$ repeats forever in base $10$. Suppose a payment service adds ten thousand micro-fees of $\$0.01$ stored as \texttt{double} and then compares \texttt{total == 100.00} with \texttt{==}. Explain in two sentences why binary floating point rounds $\frac{1}{100}$, and why storing integer minor units (\texttt{1} cent $= \texttt{1\_i64}$) makes addition exact.
\end{exercise}

\chapterline
